% !TEX encoding = UTF-8
% !TEX program = pdflatex
% !TEX spellcheck = pt_BR

% https://br.mirrors.cicku.me/ctan/macros/latex/contrib/europasscv/europasscv.pdf
\documentclass[portuguese,a4paper,nologo]{europasscv}

\usepackage[portuguese]{babel}
\usepackage[backend=biber,autolang=hyphen,sorting=none,style=numeric,maxbibnames=99,doi=false,isbn=false,maxcitenames=2]{biblatex}
\usepackage{csquotes}
\usepackage{europasscv-bibliography}


\ecvname{Elias De Nadai Nalesso Silva}
% \ecvaddress{Rua Maria Eleonora Pereira, 1170. Jardim da Penha, Vitória, ES - Brasil}
\ecvaddress{Jardim da Penha, Vitória, ES - Brasil}
% \ecvtelephone{(+555) 555 555}
% \ecvmobile{+55 27 99859-3556}
% \ecvworkphone{(+555) 123 456}
\ecvemail{pro.eliasdns@gmail.com}
% \ecvhomepage{https://lattes.cnpq.br/4361304327656565/}
% \ecvhomepage{https://portfolio.eliasdns.com/}
% \ecvgitpage{https://gist.github.com/Eliasdns/}
\ecvgithubpage{https://github.com/Eliasdns/}
% \ecvgitlabpage{gitlab.com/}
\ecvlinkedinpage{https://www.linkedin.com/in/eliasdns/}
\ecvim{WhatsApp}{https://wa.me/5527998593556}

\ecvgender{Masculino}
\ecvdateofbirth{15/08/1999}
% \ecvnationality{Brasileira}


\ecvfootnote{}

\begin{document}
% \setlength{\emergencystretch}{3em}

\begin{europasscv}
  \ecvpersonalinfo


  \ecvsection{EXPERIÊNCIA}

    \ecvtitle{Maio 2025 -- Fevereiro 2026}{Globalsys | Vixpar (Grupo Águia Branca) · Smartsourcing}
    \ecvitem{}{\textbf{Cargo:}
      \quad Dev Python Jr. e Estágio Dev Python
      \quad ·
      \quad Atuando para o projeto \textit{\underline{\href{https://play.google.com/store/apps/details?id=vix.projetovapp}{V1}}}
      (app de mobilidade urbana do Grupo Águia Branca)
    }

    \ecvitem{}{\begin{ecvitemize}
      \setlength{\itemsep}{0.4em}
      \setlength{\parsep}{0em}
      % \setlength{\parskip}{0.2em}

      \item Desenvolvi e evoluí um \textbf{dashboard} de \textbf{BI}
        e monitoramento de \textbf{KPIs} em \textbf{tempo real}
        para o app \textit{\underline{\href{https://play.google.com/store/apps/details?id=vix.projetovapp}{V1}}}
        % (app de mobilidade urbana do Grupo Águia Branca),
        integrando processamento, análise e visualização interativa de dados com boas práticas de desenvolvimento e implantação de aplicações modernas.

      \item Participei da construção conjunta de \textbf{relatórios}
        executivos em PDF (encadernáveis) com LaTex
        e de reuniões de \textbf{apresentação de resultados} para stakeholders e pessoas do negócio,
        incluindo \textbf{insights acionáveis}, \textbf{Data Storytelling} com slides (MS 365 PowerPoint)
        e \textbf{retórica} adequada.

      \item Realizei, diversas vezes, a \textbf{Análise Exploratória} dos Dados (EDA),
        \textbf{análise estatística},
        modelagem de \textbf{séries temporais}, \textbf{análise preditiva}
        e processamento e análise de dados \textbf{geolocalizados},
        investigando tendências, sazonalidades, padrões e anomalias.
        Utilizando SQL, Pandas, GeoPandas, \textit{\underline{\href{https://h3geo.org/}{H3}}}, NumPy, Statsmodels, Seaborn, Plotly e PyDeck.

      \item Desenvolvi diversas \textbf{visualizações interativas} dos dados,
        mapas de calor (choropleth maps) e mapas tridimensionais para o dashboard,
        usando Plotly, PyDeck e \textit{\underline{\href{https://h3geo.org/}{H3}}},
        além de visualizações tabelares interativas utilizando o \textit{\underline{\href{https://www.ag-grid.com/}{AG Grid}}}.

      \item Desenvolvi consultas em \textbf{SQL} / \textbf{T-SQL}
        (ocasionalmente com queries acima de 100 linhas)
        ao banco de dados MS \textbf{SQL Server}
        (banco com magnitude de centenas de tabelas e certas tabelas com milhões de registros e mais de uma centena de colunas),
        utilizando SQLAlchemy para as conexões e a respectiva autenticação via Microsoft Entra ID.

      \item Fiz parte do desenvolvimento e estruturação do front-end com Streamlit.

      \item Implementei \textbf{cache} multicamadas
        com \textbf{Redis} e na própria app com \textit{LRU-Cache} e \textit{Singleton};

      \item Estruturei \textbf{testes automatizados} com \textbf{Pytest},
        promovendo manutenibilidade e confiabilidade do dashboard.

      \item Implementei a \textbf{observabilidade} da aplicação com OTel (\textbf{OpenTelemetry})
        no lado do código (inclusive criando utilitários e métrica própria
        seguindo as convenções semânticas e de nomenclatura do OTel e do OpenMetrics/Prometheus)
        e integração com o \textbf{Azure Monitor} (Application Insights),
        além da implantação e uso de uma stack completa de observabilidade localmente (ambiente de dev.)
        (via Docker Compose) usando \textbf{Prometheus}, \textbf{Grafana}, Loki, Tempo (LGTM Stack)
        e Alloy (a distro de OTel Collector da Grafana Labs).

      \item Contribuí para a integração com o cluster k8s e com o pipeline de \textbf{CI/CD} no \textbf{Azure DevOps},
        debugando e corrigindo problemas do lado do container da app.

      \item Sendo provavelmente as minhas maiores contribuições ao projeto do dashboard:
        A definição da maior parte das dependências e dos padrões organizacionais e arquiteturais
        e a escrita e manutenção dos principais arquivos de configuração e implantação.
        % como \texttt{pyproject.toml}, \texttt{Dockerfile}, \texttt{docker-compose.yaml}, etc.

      \item Trabalhei também na \textbf{automação} (\textbf{RPA}) para aquisição de dados climáticos,
        usando a API do Open-Meteo, e na ingestão desses dados no \textbf{Databricks}.
        % usando a SQL Statement API.

    \end{ecvitemize}}


    \ecvtitle{Setembro 2023 -- Fevereiro 2025}{Top Serviços Digitais · Autônomo}
    \ecvitem{}{\textbf{Cargo:}\quad Python Fullstack Developer}

    \ecvitem{}{\begin{ecvitemize}
      \setlength{\itemsep}{0.3em}

      \item Trabalhei sob a metodologia \textbf{Scrum} e criei um sistema web de ponta-a-ponta com \textbf{Django},
        \textbf{Bootstrap}, Vanilla \textbf{JavaScript} e infra na GCP (\textbf{Google Cloud Platform}),
        que faz o controle de clientes, usuários, contratos, cobranças e pagamentos.
        Usando \textbf{UML} p/ modelagem das entidade, do banco relacional e dos casos de uso,
        prototipação da UI com \textbf{Figma}, \textbf{Pytest} p/ testes unitários, Mypy p/ tipagem estática, etc.

      \item Também auxiliei em um projeto de aplicativo mobile com \textbf{React Native}
        e back-end em \textbf{Express.js} criando sobretudo soluções em DevOps e infraestrutura.

      % {\ecvhighlight{Tecnologias}: }
    \end{ecvitemize}}


    \ecvtitle{Janeiro 2017 -- Setembro 2023}{Ilha Informática · Freelancer}
    \ecvitem{}{\textbf{Cargo:}\quad Analista de TI e Técnico em manutenção de hardware}

    \ecvitem{}{\begin{ecvitemize}
      \setlength{\itemsep}{0.3em}

      \item Criei um site de ponta-a-ponta utilizando \textbf{Django},
        \textbf{Bootstrap} e \textbf{jQuery} que automatiza parte dos processos da empresa,
        faz parsing do texto exportado pelo software usado p/ detecção de hardware (CPU-Z) com \textbf{RegEx}
        e integra essa solução com o sistema web de gestão (ERP) via API REST.

      \item Trabalhei no provisionamento de um cluster de servidores \textbf{CentOS} 7 voltado à HPC, e a sua infraestrutura.

      \item Na implantação e administração de \textbf{firewall} \textbf{pfSense}.

      \item Na administração de servidores \textbf{Windows Server} 2016 e \textbf{Active Directory}.

      \item Em diversas soluções de \textbf{virtualização}, à exemplo, Physical to virtual (P2V) com VirtualBox.

      \item Na implantação e manutenção de soluções de \textbf{CFTV} Hikvision e Intelbras.

      \item Em \textbf{Help desk} Level 1 e 2.

      \item Na manutenção de um site que utiliza \textbf{PHP}, \textbf{MySQL} e \textbf{Vanilla JavaScript}.

      \item Na \textbf{manutenção} de diversos desktops, servidores, macs, notebooks e eventualmente de smartphones e tablets,
        na solução de diversos tipos de problemas de hardware e de software,
        na implantação dos SOs e softwares necessários,
        e na \textbf{automação} de diversos desses processos, passando à exemplo,
        pela personalização de detalhes da UI do Win10,
        via scripts \textbf{Batch}, de Registro do Windows (.reg) e \textbf{PowerShell}.

      % {\ecvhighlight{Tecnologias}: }
    \end{ecvitemize}}


  \ecvsection{HABILIDADES}

    \ecvblueitem{Soft-skills}{
      \textbf{Autodidatismo}, Aprendizado rápido,
      Inteligência emocional, Inventividade, Pensamento analítico,
      Alto nível de organização e Comunicativo.
    }


    \ecvblueitem{Desenvolvimento Web, Infraestrutura em Núvem \& DevOps}{
      Sólida experiência na criação de soluções Web de ponta-a-ponta (principalmente com \textbf{Django});
      Na criação de pipelines de \textbf{CI/CD} (com \textbf{GitHub Actions});
      Na garantia de saúde e manutenibilidade do código
        com \textbf{testes automáticos} (Pytest, nox e Coverage.py),
        linting (Ruff) e
        análise estática do código com Mypy (checagem de tipos) e Bandit (checagem de segurança);
      No desenho de soluções e arquitetura usando, por exemplo, \textbf{UML} e Figma;
      E na implementação de \textbf{observabilidade} com \textbf{OpenTelemetry}, \textbf{LGTM Stack} e Azure Application Insights.
    }
    \ecvitem{}{\ecvhighlight{Tecnologias:}
      Python 3 (Avançado);
      Django framework (Avançado) e FastAPI (Intermediário);
      CSS3 (Intermediário), Bootstrap (Intermediário) e jQuery (Intermediário);
      JavaScript (Avançado), TypeScript (Básico) e Node.js (Básico);
      Elixir (Básico) e Phoenix framework (Básico).

      Git (Avançado);
      Docker (Avançado), Docker Compose (Avançado) e Docker Swarm (Básico);
      AWS (Básico), GCP (Básico) e Oracle Cloud Infrastructure (Básico);
      GitHub Actions (Avançado);
      OpenTelemetry (Avançado), Grafana (Intermediário), Prometheus (Intermediário) e Azure Monitor (Básico).
    }


    \ecvblueitem{Manutenção de computadores, Infraestrutura On-premise \& Redes}{
      Sólida experiência em \textbf{manutenção de computadores} (sobretudo Desktops e Notebooks),
      em instalações de infraestrutura de TI,
      na \textbf{administração de servidores} (Sobretudo Ubuntu, CentOS e Windows Server)
      e na administração de \textbf{firewall}s (pfSense);
    }
    \ecvitem{}{\ecvhighlight{Tecnologias:}
      Linux (Avançado) e Servidores Ubuntu e CentOS (Avançado);
      Windows (Avançado) e Windows Server (Básico);
      Firewall pfSense (Intermediário);
      Instalações de CFTV (Básico);
      Montagem e Manutenção de Computadores (Avançado).
    }


    \ecvblueitem{Automação}{
      Sólido domínio dos sistemas Linux e Windows; Sólida experiência na criação de \textbf{scripts} p/ automação;
      Forte domínio de \textbf{Python} e suas bibliotecas;
      Sólida experiência em processamento de texto com \textbf{RegEx} (incluindo Shell Parameter Expansion e tools como sed e grep).
      Conhecimento para criação e consumo de \textbf{APIs},
        para \textbf{automação de browsers}
        e em \textbf{OCR} (Reconhecimento Óptico de Caracteres).
    }
    \ecvitem{}{\ecvhighlight{Tecnologias:}
      Shell script (Avançado), Batch script (Avançado) e Powershell (Intermediário);
      Selenium (Básico) e Playwright (Intermediário);
      OCR (Reconhecimento Óptico de Caracteres) com Tesseract (Intermediário) e Google Cloud Vision API (Intermediário);
    }


    \ecvblueitem{Análise de Dados}{
      Sólida experiência
      no desenvolvimento e implantação de \textbf{Dashboards interativos},
      em \textbf{Data Visualization},
      no apoio e construção de \textbf{relatórios analíticos},
      na criação de \textbf{Data Storytelling}
      e na \textbf{comunicação} com stakeholders técnicos e de negócio (e \textbf{apresentação de resultados});
      Com experiência prática em \textbf{Análise Exploratória} de Dados (EDA),
      \textbf{Análise Estatística},
      Modelagem de \textbf{Séries Temporais}, \textbf{Análise Preditiva}
      e processamento e análise de dados \textbf{Geolocalizados},
      investigando tendências, sazonalidades, padrões e anomalias
      e gerando \textbf{insights acionáveis} para o negócio.
    }
    \ecvitem{}{\ecvhighlight{Tecnologias:}
      SQL / T-SQL (Intermediário) e MongoDB (Intermediário);
      Redis (Intermediário) e Elasticsearch (Básico);
      Pandas (Avançado) e Geoprocessamento c/ GeoPandas (Intermediário);
      Power BI (Básico);
      Plotly (Intermediário), PyDeck (Intermediário)
        e \textit{\underline{\href{https://www.ag-grid.com/}{AG Grid}}} (Intermediário);
      Front-end com Streamlit.
    }


    \ecvblueitem{Outras}{
      Prototipação e Desenvolvimento de soluções \textbf{eletrônicas} e de \textbf{IoT} com Arduino;
      \textbf{Pentest} e \textbf{hacking} de vulnerabilidades comuns, como SQL Inject, XSS, CSRF, SSRF, etc
      e uso de ferramentas como NMap, Wireshark, tcpdump, Metasploit Framework, SQLMap, Burp Suite, etc;
      Criação de planilhas (\textbf{Excel}) envolvendo gráficos e entre outros elementos mais avançados.
    }
    \ecvitem{}{\ecvhighlight{Tecnologias:}
      C (Básico);
      Arduino (Básico);
      MS Excel (Intermediário);
      Ferramentas para Pentest e Cibersegurança (Intermediário).
    }


  \ecvsection{FORMAÇÃO ACADÊMICA}
    \ecvtitle{2024 -- Presente}{Bacharel em Ciência da Computação -- FAESA}
    % \ecvitem{}{\textbf{TCC:} }
    \ecvitem{}{\ecvhighlight{Destaques:}
      2 participações no \href{https://hackfaesa.com.br/}{\textit{\underline{Hack FAESA}}} (maior hackathon universitário do estado),
      % TODO: link "oficial" acerca do projeto GEMA
      participação no projeto de extensão GEMA,
      membro da diretoria da atlética de e-Sports (e co-capitão do time de CS2)
      e coeficiente de rendimento (CR) de 9,52 ao longo do primeiro ano do curso.
    }

    \ecvtitle{2015 -- 2017}{Ensino Médio -- IFES Campus Vitória}
    \ecvitem{}{Ensino médio Integrado ao Técnico em Eletrotécnica (Incompleto)}
    \ecvitem{}{\ecvhighlight{Destaques:}
      Participação em diversas olimpíadas científicas nacionais, como OBMEP (matemática), OBA (astronomia) e OBI (informática);
      Participação como ouvinte das aulas de Cálculo I da turma de Bacharel em Engenharia do IFES no turno vespertino;
      Co-criação e parte de um grupo de estudos em informática.
    }


  \ecvsection{DIPLOMAS \& CERTIFICADOS}
    \ecvtitle{Abril 2024}{\small Certificado -- LINUXtips -- Mês do Kubernetes 2024}
    \ecvitem{}{Treinamento sobre Kubernetes e o enorme ecossistema de ferramentas para um (ambiente) cluster k8s.}

    \vspace{-0.8em}
    \ecvtitle{Janeiro 2023}{\small Curso -- LINUXtips -- Descomplicando o Docker}
    \ecvitem{}{Treinamento aprofundadado sobre Docker, Docker Compose e Docker Swarm.}

    \vspace{-0.8em}
    \ecvtitle{Abril 2022}{\small \href{http://ude.my/UC-af182aa6-e88f-490b-87da-511e75a52974}{
      Curso -- Udemy -- Elixir e Phoenix
    }}
    \ecvitem{}{Curso para criação de uma aplicação web completas com a linguagem Elixir e o framework Phoenix, já com autenticação e testes automáticos.}

    \vspace{-0.8em}
    \ecvtitle{Abril 2022}{\small \href{http://ude.my/UC-4ace7e03-86d0-44c3-abcb-ee1dfb6dbc6f}{
      Curso -- Udemy -- Programação Concorrente e Assíncrona com Python
    }}
    \ecvitem{}{Formação sobre concorrência, paralelismo e diversas técnicas e ferramentas para execução de código concorrente e assíncrono com Python.}


  % \ecvsection{IDIOMAS}
  %   % Quadro Europeu Comum de Referência para Línguas (QECR) (CEFR em inglês)
  %   \ecvmothertongue{Português}
  %   \ecvlanguageheader
  %   \ecvlanguage{Inglês}{A2}{B2}{A1}{A1}{A2}
  %   % \ecvlanguagecertificate{TOEFL -- Score: 89}
  %   \ecvlastlanguage{Espanhol}{B1}{B1}{A2}{A2}{A1}
  %   \ecvlanguagefooter


\end{europasscv}
\end{document}
